\chapter{Introduction}
\label{sec:einleitung}

Outdoor cats occasionally bring prey animals, such as mice or birds, into the house, which can be very unpleasant for the owner. These prey animals are often difficult to recognise in camera footage because they are small or occluded. When one tries to solve this problem with pure object detection, individual frames alone often prove to be error-prone \cite{jiao2022video}. Motion blur or occlusions mean that the prey cannot be recognised with the required accuracy. A snapshot is not sufficient here, which is why an analysis across several frames is required.

This bachelor thesis deals with the detection of prey animals in a cat's mouth supported by artificial intelligence. The goal of this project is to offer an alternative to conventional cat flaps and to spare the cat owner the annoying task of catching prey in the living space.

Since image recognition is, among other things, the centrepiece of this project, I have worked intensively with various algorithms. In the first step, the object detection model YOLO26 \cite{jocher2024yolo} is used. YOLO26 achieves a high accuracy while also remaining real-time capable on edge devices. Subsequently, the ByteTrack \cite{zhang2022bytetrack} algorithm is used. This algorithm tracks the detected bounding boxes over time and ensures that the system does not lose sight of the cat, even during brief occlusions. ByteTrack also ensures that each cat is uniquely identified should several cats be in the video frame. The detected regions are then evaluated by a further artificial intelligence (YOLO segmentation), which segments the prey in the image and derives a confidence score for prey from this. To increase reliability and to avoid false alarms, a decision is not made for every video frame; instead, data are collected over several video frames and a decision is made for these data as a whole. For this purpose, a custom state machine was integrated. It evaluates the history of the frames over a sliding time window and then decides whether or not the cat is carrying prey in its mouth. This thesis evaluates to what extent this combination of object detection, segmentation, object tracking and a state machine is a better solution than evaluation based purely on individual frames.
